\section{Results}
\label{sec:results}

\subsection{Audit findings}
A repository-wide search for \texttt{virtual}, \texttt{: public}, and
\texttt{template<} returned zero matches in project code, confirming no
inheritance, polymorphism, or generic code existed anywhere in the
pre-refactor firmware. Five of eleven source modules under
\texttt{AD5941\_25/src/} used \texttt{class} only as a grouping syntax
indistinguishable from a namespace, with all data members public. Two
modules were, on call-site tracing, entirely unreachable at runtime:
\texttt{command\_processing.cpp/.h} (a free-function command parser,
extending past line 700 of grammar-handling code alone) and
\texttt{electrochemical\_methods.cpp/.h} (a free-function reimplementation
of CA/SWV/DPV bound to its own separate global-variable set, several hundred
additional lines). Both duplicated functionality that a class-based
implementation already provided on the live path
(\texttt{C\_Communication} and \texttt{C\_CA}/\texttt{C\_DPV}/\texttt{C\_SWV}
respectively) against a disjoint set of global state -- a correctness hazard
as well as a maintenance one, since the two copies could silently diverge on
a nominally shared parameter. The AFE configuration/sampling/conversion
routine (\texttt{ConfigDCMeasurement}, \texttt{MeasureCurrentRaw},
\texttt{RawToCurrent}) was found byte-for-byte identical, modulo
class-qualification and comments, across \texttt{C\_CA}, \texttt{C\_DPV},
and \texttt{C\_SWV} by direct line-level diff. \texttt{Interface\_SetLed()}
was confirmed, by inspection, to ignore its \texttt{color} argument
unconditionally.

\subsection{Refactor outcome}
Table~\ref{tab:refactor} summarizes the change. The shared-base-class
extraction removed two of the three duplicate copies of the
$\sim$200-line AFE routine, leaving one. Dead-code removal deleted
\texttt{command\_processing.cpp/.h} and
\texttt{electrochemical\_methods.cpp/.h} in full. The refactored firmware
compiles cleanly against the same toolchain and board target
(\texttt{rp2040:\allowbreak rp2040:\allowbreak seeed\_\allowbreak
xiao\_\allowbreak rp2040}) with \textbf{960 fewer bytes}
of compiled flash (124\,360 vs.\ 125\,320 bytes; RAM usage unchanged at
170\,748 bytes, as expected since no data layout changed).

\begin{table}[h]
\centering
\footnotesize
\setlength{\tabcolsep}{4pt}
\caption{Refactor summary.}
\label{tab:refactor}
\begin{tabular}{@{}p{4.3cm}p{1.7cm}p{1.9cm}@{}}
\toprule
Change & Before & After \\
\midrule
Copies of AFE config/sample/convert routine & 3 (verbatim) & 1 (shared base class) \\
Dead free-function modules (command parser, electrochemical methods) & 2 & 0 \\
Status-LED color argument & Ignored (always white) & Honored (10-color table) \\
Compiled flash size & 125\,320 B & 124\,360 B ($-960$ B) \\
Compiled RAM usage & 170\,748 B & 170\,748 B (unchanged) \\
\bottomrule
\end{tabular}
\end{table}

\subsection{Why the dummy-cell traces are straight, not peaked}
Before presenting the four-method comparison, one visual point needs stating
explicitly because it looks, at first glance, like an artifact: the CV
trace in Figure~\ref{fig:before-after} is a straight line with no
oxidation/reduction peaks -- not the peaked, hysteretic ``duck'' shape a
cyclic voltammogram of a redox-active analyte (e.g.\ ferro/ferricyanide)
shows. This is the expected, physically correct result, not a
measurement or plotting error, and it is not specific to this instrument:
it is a property of the dummy cell, established independently of any
potentiostat. A Randles dummy cell and a plain resistor are passive
$R$/$RC$ networks with no Faradaic reaction; peaks in a CV trace arise
specifically from the potential-dependent rate of an electron-transfer
reaction, which a dummy cell does not contain by construction. Wong and
Zhou's direct theoretical treatment of CV on resistor/capacitor equivalent
circuits derives exactly this result -- a passive $R$/$RC$ network under a
linear potential sweep produces a current that is a linear (resistive
branch) or elliptical/parallelogram-like (capacitive branch) function of
the applied potential, structurally incapable of the peaked, asymmetric
shape a Faradaic reaction produces~\cite{wong2024cyclic}. Commercial dummy
cells are built, and documented, on the same premise: manufacturer
reference cells are explicitly specified to give a predictable, calculable
$R$/$RC$ response rather than a redox-like one, precisely so that the
instrument, not the analyte, is what is under test~\cite{gamry2011universal},
and independent round-robin characterization of dummy cells for
electrochemical instrumentation reports the same linear/predictable
behavior as the reference case against which instrument error is
measured~\cite{kasper2022electrochemical}. What a dummy cell can and should
show is a closed loop traced twice, once on the forward (increasing
potential) leg and once on the reverse leg, and the two panels below plot
both legs explicitly (solid vs.\ dotted/dash-dot) rather than as a single
merged trace. They overlap almost exactly, which is itself informative: it
means the RC time constant of both loads ($R_{ct}C_{dl}\approx
\SI{330}{\micro\second}$ for the Randles cell) is fast relative to the
\SI{100}{\milli\volt\per\second} scan rate used here, so hysteresis from
capacitive charging current is below the trace width at this scan rate --
consistent with, not contradicted by, the dummy cell's known $R$/$RC$
composition.

\subsection{Hardware validation: refactor vs.\ pre-refactor firmware}
Figure~\ref{fig:before-after} overlays all four methods against the Randles
dummy cell for the pre-refactor and post-refactor firmware. For CV
($n=305$ points, \SI{-200}{\milli\volt} to \SI{600}{\milli\volt},
\SI{5}{\milli\volt} step, both forward and reverse legs shown), the mean
relative difference between firmware versions is $\mathbf{0.070\,\%}$ and
the maximum is $\mathbf{2.83\,\%}$ (restricting to $|I|>1$ to avoid
division blow-up at the zero crossing; the maximum occurs only immediately
adjacent to that crossing). CA shows both firmware versions settling to the
same quantized current level ($\approx -2.95\,\mu$A) with sample-to-sample
jitter of about one ADC code in either direction -- the step response of a
near-resistive load settles within one sample period, so what is visible is
quantization noise, not a decay curve, and it is statistically
indistinguishable between firmware versions. SWV shows a sharp initial
point at the first step (a known transient at sweep start, present
identically in both versions) followed by both traces tracking each other
within the method's noise floor. DPV shows a mild local maximum near
\SI{20}{\milli\volt} in both firmware versions -- an artifact of the pulse
sequencing at the first few steps, not a redox peak, and reproduced
identically before and after the refactor. In no method does the
post-refactor trace diverge from the pre-refactor trace by more than
ordinary run-to-run noise.

\begin{figure}[h]
\centering
\includegraphics[width=\linewidth]{fig1_all_methods_before_after}
\caption{All four methods against the Randles dummy cell, pre- vs.\
post-refactor firmware, same physical instrument and load. CV shows both
sweep legs explicitly (forward: solid/dashed; reverse: dotted/dash-dot).}
\label{fig:before-after}
\end{figure}

\subsection{Corroboration on a second, independent load}
Figure~\ref{fig:two-loads} compares the refactored firmware's response on
the Randles dummy cell ($R_s+R_{ct}\approx\SI{10.56}{\kilo\ohm}$ effective
DC resistance) against a plain $\SI{10}{\kilo\ohm}$ resistor substituted in
its place, across all four methods. For CV, the two loads differ in
nominal DC resistance by a factor of $10560/10000 = 1.056$; a linear
regression of reported current against applied potential gives slopes of
$0.09422$ (Randles) and $0.09986$ (resistor) current-units per mV, a ratio
of $0.9435$ -- within $0.9\,\%$ of the $1/1.056 = 0.9470$ value Ohm's law
predicts from the two loads' nominal resistance ratio alone. CA again
settles to the same quantized level for both loads, as expected since both
are within $6\,\%$ of each other in resistance and CA's fast RC settling
makes it largely insensitive to that difference at this current range. SWV
and DPV track each other in shape across both loads, with point-to-point
differences of the same magnitude as the before/after comparison in
Figure~\ref{fig:before-after}. This is not a precision calibration claim
(that is the scope of the companion manuscript~\cite{clement2026linearity});
it is a sanity check that the refactored firmware responds correctly, in
the correct direction and roughly the correct magnitude, to a real load it
was never exercised against during development of the refactor.

\begin{figure}[h]
\centering
\includegraphics[width=\linewidth]{fig2_all_methods_two_loads}
\caption{Refactored firmware only: response to two independent real loads
(Randles dummy cell vs.\ plain \SI{10}{\kilo\ohm} resistor), all four
methods.}
\label{fig:two-loads}
\end{figure}
